\section{Finite clusters and their positive charge}\label{f:cluster-section}

The objects in this section are finite clusters, rather than a product
over all zeros of an entire function.  This distinction makes the
boundary assumptions explicit and retains the influence of the exterior
zeros through a holomorphic unit.  The construction consolidates
\cite{note04}.

\begin{definition}[Proper spatial cluster]\label{f:proper}
Let $D\subset\mathbb C$ be a bounded disk or rectangle, and let
$\Omega_R=\{z:|z-a|<R\}$.  Suppose $H$ is holomorphic on a neighborhood
of $\overline D\times\overline\Omega_R$ and
\[
 H(t,z)\ne0\qquad
 (t\in\overline D,\ |z-a|=R).
\]
The zeros of $H(t,\cdot)$ in $\Omega_R$, counted with multiplicity,
form a proper cluster.  Their number $m$ is independent of $t$.
A collar of radius $r<R$ means that all these zeros lie in
$\Omega_r=\{|z-a|<r\}$ throughout $\overline D$.
\end{definition}

The count is the integer-valued continuous function
$(2\pi i)^{-1}\oint H_z/H\,dz$.  In particular, two counts, one on each
of the circles of radii $r$ and $R$, certify a collar; an inner count
alone does not exclude additional roots in the annulus.

\begin{lemma}[Preparation without root labels]\label{f:preparation}\label{f:cluster}
A proper cluster has a monic polynomial
\[
 P(t,z)=z^m+p_1(t)z^{m-1}+\cdots+p_m(t)
\]
with coefficients holomorphic in $t$, and
$H(t,z)=U(t,z)P(t,z)$, where $U$ is holomorphic and nonvanishing on
the product.  The polynomial is uniquely determined by the cluster.
\end{lemma}

\begin{proof}
For $k\ge1$ the contour integrals
\begin{equation}\label{f:root-moments}
 S_k(t)=\frac1{2\pi i}\oint_{|\zeta-a|=R}
       \zeta^k\frac{H_z(t,\zeta)}{H(t,\zeta)}\,d\zeta
\end{equation}
are holomorphic and equal the sums of the $k$th powers of the roots,
including multiplicity.  Newton's identities determine the elementary
symmetric functions recursively from $S_1,\ldots,S_m$, hence construct
$P$ with holomorphic coefficients.  On each fixed time fiber $H/P$
has removable singularities at exactly the selected roots and has no
zeros in the disk.  Joint holomorphy follows, without choosing those
roots, from the Cauchy representation
\[
 U(t,z)=\frac1{2\pi i}\oint_{|\zeta-a|=R}
        \frac{H(t,\zeta)}{P(t,\zeta)(\zeta-z)}\,d\zeta,
       \qquad |z-a|<R.
\]
For each $t$ this is the removable quotient, and the integrand is
holomorphic in its parameters.  Near the spatial wall the quotient
itself supplies the extension.  Compactness and the strict boundary
condition allow the same construction on a slightly enlarged time
neighborhood.  Uniqueness follows from the multiset of roots.
\end{proof}

Define the finite discriminant by
\begin{equation}\label{f:discriminant-definition}
 \Delta_P(t)=\operatorname{Disc}_z P(t,z)
    =\prod_{i<j}(z_i(t)-z_j(t))^2.
\end{equation}
The symmetric expression is holomorphic even when the roots cannot be
labeled holomorphically.  We set $\Delta_P=1$ if $m\le1$.

\begin{theorem}[Boundary charge and discriminant winding]
\label{f:charge-winding}\label{f:charge}
In addition to properness, suppose that $(H,H_z)$ has no zero on
$\partial D\times\overline\Omega_R$.  Give complex two-space its
complex orientation and the product boundary its induced orientation.
The degree of
\[
 \frac{(H,H_z)}{\sqrt{|H|^2+|H_z|^2}}:
       \partial(D\times\Omega_R)\longrightarrow S^3
\]
is well defined and equals
\begin{equation}\label{f:winding-formula}
 Q_{D\times\Omega_R}
   =\frac1{2\pi i}\oint_{\partial D}
              \frac{\Delta_P'(t)}{\Delta_P(t)}\,dt.
\end{equation}
It is the sum of the local complex intersection multiplicities of
$(H,H_z)$, and is therefore a nonnegative integer.
\end{theorem}

\begin{proof}
The side wall contains no zero of $H$, and the remaining boundary
contains no joint zero by hypothesis.  The discriminant is nonzero
on $\partial D$.  It is consequently not identically zero and has
only finitely many zeros in $\overline D$.  Above each such time there
are only finitely many cluster roots, so the joint zeros are isolated.

Fix one such time $t_0$ and put
$\mathcal R=\mathbb C\{t-t_0\}$, the ring of convergent power series.
The module $\mathcal A=\mathcal R[z]/(P)$ is free of rank $m$ over
$\mathcal R$.  Multiplication by $P_z$ on $\mathcal A$ has determinant
\[
 \operatorname{Res}_z(P,P_z)
    =(-1)^{m(m-1)/2}\Delta_P.
\]
Indeed, at a time with distinct roots multiplication is diagonal in
the evaluation coordinates, with determinant $\prod_iP_z(z_i)$;
the identity then extends holomorphically.  The ring $\mathcal R$ is
a discrete valuation ring.  Elementary row and column reduction puts
any matrix with nonzero determinant into diagonal form with entries
units times powers $(t-t_0)^{e_j}$.  Thus the order of its determinant
is the sum of the $e_j$, which is also the length of its cokernel.
Applying this to multiplication by $P_z$ gives
\[
 \operatorname{ord}_{t_0}\Delta_P
       =\operatorname{length}_{\mathcal R}\mathcal A/(P_z).
\]
Small disjoint spatial disks around the distinct roots at $t_0$
split this quotient into its local factors.  In each factor, analytic
division by the monic local polynomial identifies its length with
\[
 \dim_{\mathbb C}
 \frac{\mathbb C\{t-t_0,z-b\}}{(P,P_z)}.
\]
Since $U$ is a unit, $(H,H_z)=(UP,UP_z+U_zP)$ and $(P,P_z)$
generate the same ideal.  These dimensions are therefore precisely
the local intersection multiplicities of the original jet.

For completeness, the topological interpretation uses the usual local
degree property of a holomorphic square map: at an isolated zero its
degree is the dimension of its local analytic quotient.  One way to
see the sign and multiplicity is to perturb the target to a nearby
regular value in an isolating neighborhood.  The local intersection
number counts the preimages with their multiplicities; at a regular
preimage the real Jacobian determinant is
$|\det_{\mathbb C}DF|^2>0$.  Conservation of the intersection number
and homotopy invariance of boundary degree give the asserted equality
at the original value.  Additivity over disjoint isolating neighborhoods
then identifies the product-boundary degree with the sum of the above
lengths.  The scalar argument principle sums
$\operatorname{ord}_{t_0}\Delta_P$ and proves
\eqref{f:winding-formula}.  Corners of a product rectangle may be
rounded within a boundary neighborhood free of joint zeros, without
changing either degree.
\end{proof}

The heat equation determines the local multiplicity in this formula.
Here and below the time normalization is $H_t=-H_{zz}$.

\begin{proposition}[Hermite birth and local charge]\label{f:hermite}
Suppose $H(t_0,b+w)=c_mw^m+O(w^{m+1})$, where $c_m\ne0$ and
$m\ge2$.  The event $(t_0,b)$ is an isolated zero of $(H,H_z)$.
For its local $m$-root cluster, writing $s=t-t_0$ and
$N=m(m-1)/2$, one has
\begin{equation}\label{f:hermite-discriminant}
 \Delta_P(t_0+s)=d_m s^N+O(s^{N+1}),\qquad
 d_m=2^N\prod_{i<j}(\lambda_i-\lambda_j)^2>0,
\end{equation}
where $\lambda_1,\ldots,\lambda_m$ are the roots of the probabilists'
Hermite polynomial $\operatorname{He}_m$.  The local jet charge is $N$.
\end{proposition}

\begin{proof}
The heat equation recursively determines the Taylor coefficients in
$s$ from the initial spatial Taylor series.  Terms arising from $w^\ell$
have weighted degree $\ell$ when $w$ has weight one and $s$ has weight
two.  The leading weighted polynomial is
\[
 c_m e^{-s\partial_w^2}w^m
   =c_m(2s)^{m/2}\operatorname{He}_m
                  \left(\frac{w}{\sqrt{2s}}\right).
\]
This is a polynomial identity, so the notation involving square roots
introduces no choice into the polynomial itself.

The polynomial $\operatorname{He}_m$ has $m$ distinct real roots.
Indeed, its Rodrigues formula and integration by parts imply
orthogonality to every polynomial of degree less than $m$ for the
positive weight $e^{-x^2/2}$.  If it had fewer than $m$ distinct real
roots of odd multiplicity, multiply it by the product of the linear
factors at just those roots.  The resulting real polynomial would
have a constant nonzero sign except at finitely many points, so its
weighted integral could not vanish.  Orthogonality contradicts this.
Degree $m$ now forces all $m$ roots to be real and simple.

Set $s=\varepsilon^2$ and $w=\varepsilon x$.  Taylor convergence gives,
uniformly for bounded $x$,
\[
 \varepsilon^{-m}H(t_0+\varepsilon^2,b+\varepsilon x)
 =c_m2^{m/2}\operatorname{He}_m(x/\sqrt2)+O(\varepsilon).
\]
The left side extends holomorphically to $\varepsilon=0$.
The simple-root implicit function theorem gives all $m$ cluster roots
as
$w_i=\varepsilon(\sqrt2\lambda_i+O(\varepsilon))$.
They are distinct for sufficiently small nonzero $\varepsilon$.
Their squared differences have product
$d_m\varepsilon^{2N}(1+O(\varepsilon))$.
The discriminant is holomorphic in $s=\varepsilon^2$, so its Taylor
series has only even powers of $\varepsilon$; this proves
\eqref{f:hermite-discriminant}.  It also proves local isolation of
the event.  Theorem~\ref{f:charge-winding} identifies its charge with
the order $N$.
\end{proof}

Consequently, on any product satisfying Theorem~\ref{f:charge-winding}
for a heat family,
\begin{equation}\label{f:positive-sum}
 Q_{D\times\Omega_R}
    =\sum_{(t,b):\,H(t,b)=H_z(t,b)=0}\binom{m(t,b)}2.
\end{equation}
Interior event charges do not cancel.  Reality and spatial parity do
not change this: simultaneous conjugation in two complex coordinates
has real orientation sign $+1$, and spatial sign reversal is complex
linear, also of sign $+1$.  These symmetries can pair equal charges,
but cannot force their sum to vanish.

\section{The two-root equation and the real-cap theorem}
\label{f:two-root-section}

Assume a proper cluster consists of two roots in a collar
$0<r<R$.  Write $H=UP$, $P=(z-z_1)(z-z_2)$, and
$\Delta=(z_1-z_2)^2$.  Set $q=U_z/U$ and, on the outer circle,
$f=H_z/H$.  The exterior unit $U$ is retained throughout; replacing
it by a constant would change the equation for the original cluster.

\begin{proposition}[Regular affine discriminant flow]\label{f:affine}
The holomorphic coefficient
\begin{equation}\label{f:A-contour}
 A(t)=\frac4{2\pi i}\oint_{|\zeta-a|=R}
        \frac{H_z(t,\zeta)}{H(t,\zeta)P(t,\zeta)}\,d\zeta
\end{equation}
is regular through collisions and satisfies
\begin{equation}\label{f:affine-equation}
                    \Delta'=8+A\Delta.
\end{equation}
For any holomorphic scalar $\kappa(t)$, $H_z$ may be replaced in
\eqref{f:A-contour} by $H_z-\kappa H$.  If on the outer wall
$|H|\ge\delta>0$ and $|H_z-\kappa H|\le L_\kappa$, then
\begin{equation}\label{f:A-bound}
          |A|\le K:=\frac{4RL_\kappa}{\delta(R-r)^2}.
\end{equation}
Every zero of $\Delta$ is simple and has derivative $8$.
\end{proposition}

\begin{proof}
At a simple root the heat equation gives $z_i'=H_{zz}/H_z$.
Differentiating $UP$ at that root yields
\[
 z_1'=\frac2{z_1-z_2}+2q(t,z_1),\qquad
 z_2'=\frac2{z_2-z_1}+2q(t,z_2).
\]
Thus
\[
 \Delta'=8+4\Delta\frac{q(t,z_1)-q(t,z_2)}{z_1-z_2}.
\]
By residues the divided difference is
$(2\pi i)^{-1}\oint q/P\,d\zeta$.  Since $f=q+P_z/P$ and
$P_z/P^2=-\partial_z(1/P)$, the integral is unchanged if $q$ is
replaced by $f$.  This proves the equation away from events.
The discriminant is not identically zero by
Proposition~\ref{f:hermite}; alternatively, a persistent double
$H=U(z-c(t))^2$ would give $H_t(t,c(t))=0$ and
$H_{zz}(t,c(t))=2U(t,c(t))\ne0$, contradicting heat flow.
Both sides of \eqref{f:affine-equation} are holomorphic, so the equation
extends through its isolated events.

The scalar subtraction is exact because $\oint1/P\,d\zeta=0$:
the rational function $1/P$ has no exterior finite poles and decays
quadratically at infinity.  On the outer wall,
$|P|\ge(R-r)^2$.  Estimating the contour length $2\pi R$ gives
\eqref{f:A-bound}.  Substitution of $\Delta=0$ in the equation gives
$\Delta'=8$.
\end{proof}

There is no inverse internal root separation in
\eqref{f:A-contour} or \eqref{f:A-bound}.  The bound does pay for
the exterior collar and the spatial-wall floor.  A scalar centering
$\kappa$ improves it only if the corresponding coupled numerator is
actually estimated more sharply.

\begin{theorem}[Reduction of complex charge to two real cap signs]
\label{f:real-cap}
Let
\[
 D=\{x+iy:t_-<x<t_+,\ |y|<\tau\},\qquad t_-<t_+,
\]
and suppose the proper two-root and collar hypotheses hold on a
neighborhood of its closure.  Let $a\in\mathbb R$ and assume reality
$H(\overline t,\overline z)=\overline{H(t,z)}$.  It suffices that
\begin{equation}\label{f:phase}
 \sup_{x+iy\in\overline D}
   \left|\int_0^y\operatorname{Re} A(x+is)\,ds\right|<\frac\pi2.
\end{equation}
In particular, $|A|\le K$ and $K\tau<\pi/2$ suffice.  Then there
is at most one cluster event in $\overline D$, and any such event has
real time and real spatial position.  If
$\Delta(t_-)\Delta(t_+)\ne0$, the entire product boundary is free
of joint zeros, and
\begin{equation}\label{f:cap-charge}
 Q_{D\times\Omega_R}
  =\frac{\operatorname{sgn}\Delta(t_+)
               -\operatorname{sgn}\Delta(t_-)}2\in\{0,1\}.
\end{equation}
\end{theorem}

\begin{proof}
Choose $t_*\in(t_-,t_+)$ and define the holomorphic primitives
\[
 J(t)=\int_{t_*}^t A(v)\,dv,\qquad
 \mu(t)=e^{-J(t)},\qquad \Psi(t)=\int_{t_*}^t\mu(v)\,dv.
\]
They are single valued on the rectangle.  Reality and uniqueness of
the prepared polynomial imply reality of $A$ on the real interval.
Integrating first horizontally and then vertically gives
\[
 \operatorname{Im}J(x+iy)=\int_0^y\operatorname{Re}A(x+is)\,ds.
\]
Therefore \eqref{f:phase} implies $\operatorname{Re}\mu>0$ on the
closed rectangle.  Convexity gives, for distinct $s,t$ in its closure,
\[
 \frac{\Psi(t)-\Psi(s)}{t-s}
      =\int_0^1\mu(s+v(t-s))\,dv,
\]
whose real part is positive.  Thus $\Psi$ is injective, including
on the boundary.  Multiplying \eqref{f:affine-equation} by $\mu$
and integrating yields
\begin{equation}\label{f:normal-form}
                 \mu(t)\Delta(t)=\Delta(t_*)+8\Psi(t).
\end{equation}
There is at most one zero.  A nonreal zero would have a distinct
conjugate zero, so any zero is real.  At a real event the real monic
quadratic has a repeated root, which is necessarily real as well.

The only possible event times on the boundary are now $t_-$ and
$t_+$, and the cap assumptions exclude them.  On real time $\mu>0$
and $(\mu\Delta)'=8\mu>0$.  Equal nonzero cap signs give no event;
negative then positive signs give exactly one.  The reversed sign
pattern cannot occur.  Each event contributes one by
Proposition~\ref{f:affine} and Theorem~\ref{f:charge-winding}, proving
\eqref{f:cap-charge}.
\end{proof}

The theorem proves an upper bound of one event, not zero events.  The
heat family $H=z^2-2t$ has $A=0$ and $\Delta=8t$; on
$|\operatorname{Re}t|,|\operatorname{Im}t|\le1/4$ its roots lie in
$|z|<1$, so $r=1,R=2$ supply a proper collar and charge one.  The
phase condition is entirely compatible with this collision.  This
distinction is central to the endpoint problem \cite{note05}.

\subsection{A complete-source strip criterion}

For the complete positive source, use the moments
\[
 M_j(T,s)=\int_0^\infty\Phi(u)u^j e^{Tu^2}\cosh(su)\,du.
\]
They are finite, and no source modes or exterior tail are omitted.
Suppose $d_r,d_R>0$ are established bounds
\[
 |H(x,z)|\ge d_s\quad
 (t_-\le x\le t_+,\ |z-a|=s),\qquad s=r,R,
\]
and a real anchor has root count two on both circles.  Put
$B_s=M_2(t_+,s)$ and $L=M_1(t_+,R)$.  Then a sufficient width is
\begin{equation}\label{f:source-strip}
 0<\tau<\min\left\{
    \frac{d_r}{B_r},
    \frac{d_R}{B_R+8RL/(\pi(R-r)^2)}\right\}.
\end{equation}
Indeed, $|e^{iyu^2}-1|\le|y|u^2$ and
$|\cos(zu)|,|\sin(zu)|\le\cosh(su)$ on a circle centered at a
real $a$ give
\[
 |H(x+iy,z)-H(x,z)|\le\tau B_s,\qquad
 |H_z(x+iy,z)|\le M_1(t_+,s).
\]
Both walls therefore have positive floors $d_s-\tau B_s$, and
continuity of their counts gives the proper two-root collar throughout
the strip.  Applying \eqref{f:A-bound} with $\kappa=0$ and
$\delta=d_R-\tau B_R$ gives $K\tau<\pi/2$ exactly from the second
inequality in \eqref{f:source-strip}.  Purely imaginary time lifting
requires no extra factor $e^{\tau u^2}$.  The cap signs still require
separate information.  The finite verified application discussed later
uses this type of complete-source payment \cite{note06}.

\section{Explicit windows at a hypothetical double}
\label{f:window-section}

The preceding theorem has a pointwise geometric contract at every
finite real height.  This section pays it in terms of the local
second jet, including when the two roots coincide \cite{note07}.
Fix $T_0,R_0>0$ and set
\begin{equation}\label{f:B-moments}
 B_j=M_j(T_0,R_0),\qquad j=2,3.
\end{equation}
The complete source gives
$|\partial_z^jH(t,z)|\le B_j$ whenever
$\operatorname{Re}t\le T_0$ and $|\operatorname{Im}z|\le R_0$.
These bounds do not depend on $\operatorname{Re}z$.

\begin{theorem}[Quantitative double window]\label{f:double-window}
Assume, conditionally, that for a real $a$,
\[
 H(0,a)=H_z(0,a)=0,\qquad
 c=\tfrac12H_{zz}(0,a)\ne0,
 \qquad q_0=|c|.
\]
Define
\begin{equation}\label{f:window-constants}
 \begin{split}
 R&=\min\{R_0/2,q_0/B_3\},\qquad r=R/2,\\
 h&=\min\left\{T_0/2,
       \frac{q_0R^2}{16\sqrt2 B_2},\frac{R^2}{128}\right\}.
 \end{split}
\end{equation}
For every $t$ with $|\operatorname{Re}t|,|\operatorname{Im}t|\le h$,
$H(t,\cdot)$ has exactly two roots inside each of the spatial circles
of radii $r$ and $R$, and their walls satisfy
\begin{equation}\label{f:double-floors}
 |H(t,z)|>\tfrac12q_0s^2\quad(|z-a|=s,\ s=r,R).
\end{equation}
On the outer wall $|H_z|<3q_0R$.  The regular coefficient satisfies
\begin{equation}\label{f:double-phase}
          |A|\le K:=96/R^2,\qquad Kh\le3/4<\pi/2.
\end{equation}
The hypotheses extend to a neighborhood of the closed product.
\end{theorem}

\begin{proof}
For $|w|=s\le R$, spatial Taylor expansion with integral remainder
and straight-line integration in complex time give
\[
 |H(0,a+w)-cw^2|\le B_3s^3/6\le q_0s^2/6,
\]
\[
 |H(t,a+w)-H(0,a+w)|\le|t|B_2
       \le\sqrt2hB_2\le q_0R^2/16=q_0r^2/4.
\]
At $s=r,R$ their sum is at most $5q_0s^2/12<q_0s^2$.
Rouch\'e's theorem against $cw^2$ gives both counts and the stronger
common floor $7q_0s^2/12$.  At the outer circle the same calculation
gives $37q_0R^2/48$.  Counting twice excludes annular roots.

For the first derivative use its spatial remainder and
$H_{zt}=-H_{zzz}$:
\[
 \begin{split}
 |H_z(t,a+w)|
 &\le2q_0R+\tfrac12 B_3R^2+\sqrt2hB_3\\
 &\le\left(\tfrac52+\frac{\sqrt2}{128}\right)q_0R<3q_0R.
 \end{split}
\]
The last step uses $RB_3\le q_0$ and $h\le R^2/128$.
Equation~\eqref{f:A-bound}, with
$\delta=q_0R^2/2$ and $R-r=R/2$, now gives $K=96/R^2$.
All required wall margins are strict and $Kh\le3/4<\pi/2$;
continuity on the compact boundary preserves the conditions on a
small enlargement.  The corner factor $\sqrt2$ was included explicitly
in the time estimates.
\end{proof}

\begin{corollary}[The compatible cap signs]\label{f:double-caps}
In Theorem~\ref{f:double-window} define, for $x\ge0$,
\begin{equation}\label{f:LU}
 L_K(x)=\frac8K(1-e^{-Kx}),\qquad
 U_K(x)=\frac8K(e^{Kx}-1),
\end{equation}
with continuous limiting values $8x$ if $K=0$.  Then
\begin{equation}\label{f:double-cap-bounds}
 \begin{split}
 L_K(h)&\le\Delta(h)\le U_K(h),\\
 -U_K(h)&\le\Delta(-h)\le-L_K(h).
 \end{split}
\end{equation}
The resulting product has charge one, with its unique event at $(0,a)$.
\end{corollary}

\begin{proof}
At the assumed double, $\Delta(0)=0$.  On real time the affine equation
has the exact solution
\[
 \Delta(t)=8\int_0^t
       \exp\left(\int_s^t A(v)\,dv\right)\,ds.
\]
The coefficient is real and bounded by $K$, so comparison with
$e^{-K|t-s|}$ and $e^{K|t-s|}$ gives
\eqref{f:double-cap-bounds}.  In particular the cap margins are
positive, and $L_K(h)\ge8he^{-Kh}$.  Theorem~\ref{f:real-cap}
gives the charge and uniqueness.
\end{proof}

This is a geometry theorem conditional on an exact double.  It cannot
also be an exclusion theorem: its actual cap signs are those of a
collision.  Improving an absolute coefficient bound changes their
quantitative margins, but cannot turn them into equal signs.

\subsection{A finite-jet test away from an assumed event}

At any real center let
$c_0=H(0,a)$, $c_1=H_z(0,a)$, $c_2=H_{zz}(0,a)/2$, and assume
$q_0=|c_2|>0$.  Choose $0<r<R\le R_0$ and
$\sqrt2h\le T_0$.  If, for $s=r,R$,
\begin{equation}\label{f:off-event}
 d_s=q_0s^2-|c_0|-s|c_1|-B_3s^3/6-\sqrt2hB_2>0,
\end{equation}
then the same Rouch\'e proof gives two roots inside each circle and
wall floors $d_s$ on the closed time square.  Moreover
\begin{equation}\label{f:off-event-A}
 L=|c_1|+2q_0R+B_3R^2/2+\sqrt2hB_3,
 \qquad K=\frac{4RL}{d_R(R-r)^2}
\end{equation}
are valid outer derivative and coefficient bounds.  If $Kh<\pi/2$,
the real-cap theorem applies once the two cap discriminants are
enclosed away from zero.  The test does not assert an event, an
internal gap, or either cap sign.  An interval implementation needs a
lower enclosure for $q_0$ and complete upper enclosures for every
error term; uncertain jet equalities cannot be promoted to an exact
double.

\subsection{Event duration and the height quantifier}

If an event is already known to occur at real time $s<t_*$ and
$|A|\le K$ between them, put $d=t_*-s$ and
$\delta=\Delta(t_*)>0$.  The preceding integral solution gives
$L_K(d)\le\delta\le U_K(d)$.  Therefore, if
$0<\delta_-\le\delta\le\delta_+$ and $K\delta_+<8$,
\begin{equation}\label{f:duration}
 \frac{\log(1+K\delta_-/8)}K
 \le d\le
 -\frac{\log(1-K\delta_+/8)}K.
\end{equation}
For $K=0$ the interval is $[\delta_-/8,\delta_+/8]$.
These inequalities require event existence, supplied for example by
opposite cap signs; a single positive cap value does not supply it.
The scalar bounds are sharp under the information $|A|\le K$, as
constant coefficients $A=K$ and $A=-K$ attain the two extremes.
This statement does not assert that either coefficient is realized by
theta.  A strictly negative event location would follow only from
data making the resulting upper bound for $s$ negative.

The moments in \eqref{f:B-moments} work for every finite real $a$.
Thus every hypothetical exact double has a positive window from
\eqref{f:window-constants}, point by point.  No uniform lower bound
for its second jet is proved.  In fact, with
$f(u)=u^2\Phi(u)$ and $C=\|f'\|_{L^1(0,\infty)}/2<\infty$,
integration by parts gives, for $a\ne0$,
\[
 q_0(a)=\frac12\left|\int_0^\infty f(u)\cos(au)\,du\right|
         \le\frac C{|a|}.
\]
Consequently this particular prescription obeys
\begin{equation}\label{f:height-scale}
 R(a)\le\frac{C}{B_3|a|},\qquad
 h(a)\le\frac{C^3}{16\sqrt2 B_2B_3^2|a|^3}.
\end{equation}
These are upper bounds for conservatively chosen window sizes, not
necessary scales for every possible construction.  On a bounded
real interval there are only finitely many zeros of the nonzero entire
function $H(0,\cdot)$, so the windows of its exact doubles there have
positive minima if the set is nonempty.  This does not compute those
minima or identify any doubles.  An all-height exclusion may work
pointwise and need not use a common negative cap, but its signed
estimate must cover every admissible finite $a$ and the associated
shrinking parameters.  The geometric quantifier is paid here; that
signed quantifier remains open.
